We study how small networks trained to execute single-source Bellman-Ford shortest paths on random graphs with 4--8 nodes behave on graphs with 16, 32 and 64 nodes. We reimplement an MLP on the zero-padded flattened adjacency matrix and message-passing networks (MPNNs) with sum or max aggregation, each trained with or without supervision of the intermediate distance after every relaxation round, and test on two graph families: constant expected degree (sparse) and growing degree (dense). With five seeds and 1000 Adam steps per run, the expected pattern is only partly reproduced. On the registered primary metric, the mean error on sparse 64-node graphs, the sum-aggregation MPNN trained on the final distance alone is lowest and has no failing sparse seed. The max-aggregation MPNN trained on the final distance alone has the lowest error in distribution and the lowest median error at 64 nodes (medians were chosen post hoc), but two of its five seeds exceed our failure threshold there. The MLP is poor even in distribution, and its relative error at 64 nodes ($0.943$ on sparse graphs) is close to that of predicting zero. Step-wise supervision did \emph{not} help in our setup: median error at 8 and 16 nodes was higher with it for both aggregators, and the sum-aggregation MPNN with step supervision produced a non-finite error on dense 64-node graphs in four of five seeds and an error of order $10^{17}$ in the fifth. None of the registered comparisons at 64 nodes is significant (all $p>0.1$). For the two step-supervised MPNNs (three seeds each; the final-only models were not swept), running more message-passing steps at test time than $n-1$ raised the error at 16 nodes in all six models and at 64 nodes in four of them, and running half as many steps lowered it. The study is small (one algorithm, one training budget, no tuning, five seeds) and the results should be read as a cautionary baseline, not a ranking.
